\begin{theorem}[d two]
\uses{d}
\label{d_two}
\lean{ArithmeticFunction.d_two}
% at /home/runner/work/trough/trough/library/PrimeNumberTheoremAnd/IwaniecKowalskiCh1.lean:211.0-219.57
\leanok
$d_2$ is the classical divisor count function $\tau$.
\end{theorem}
\begin{proof}
\uses{zeta_mul_zeta}
\leanok
By definition, $d_k$ is the $k$-fold Dirichlet convolution of $\zeta$. When $k = 2$, this means we are taking the convolution of $\zeta$ with itself twice, which gives us $\zeta * \zeta$. From the earlier theorem, we know that $\zeta * \zeta = \tau$, where $\tau$ is the divisor count function. Therefore, $d_2 = \zeta^2 = \tau$.
\end{proof}
